Both trackers follow the dominant ridge of the forward FTLE field from the
channel entrance to the bifurcation and take the same branch past it
(Fig.~\ref{fig:ocean_ftle}), with no changes to the primitives at the
ocean operating point of Table~\ref{tab:params}. The field is divergent and
not vorticity-free, so this agreement is not built in. Along both paths,
however, the fit is strain-dominated, with $\hat{D} < 0$ on about
80\% of steps and a median $|\hat{\omega}|/2\hat{r}$ of 0.15 to 0.20 from
7 to 21~h. $D$ is then close to $s_1 s_2 = \mu^2 - r^2$, where the
divergence enters only as $\mu^2$, a tenth to a fifth of $r^2$. The paths
lie within 1.3~km of each other on average.

The $s_1$ tangent also types the structure. From 7 to 14~h it was
$\mathbf{e}_2$ on 34 of 42 steps, an attracting OECS where floating
material collects \cite{12}, although forward FTLE marks repelling
lines. From 14 to 21~h it was $\mathbf{e}_1$ on 37 of 42 steps, a
repelling OECS only where the trench is also a ridge of $s_2$, on 58\% of
such steps.

The $D$ tracker lands at 23~h, and the $s_1$ tracker is at sea at 28~h.
From 2 to 6.5~h the $s_1$ tracker crossed a
rotation-dominated patch, $\hat{D} > 0$ on 70\% of steps, where its
tangent switched 12 times in 27 steps. It made 2.1~km of net progress in
7~h against 8.4~km for the $D$ tracker, then kept pace. Mean distance to
the nearest ridge point is 1.6~km for the $D$ tracker and 2.1~km for the
$s_1$ tracker. The $s_1$
tracker's grows to 3.8~km below $34.21^{\circ}$N, reached 7~h after the
$D$ tracker once the fitted strain there had halved, as Michini et al.
report on this record \cite{26}.
